%------------------------------------------------------------------------------%
\documentclass[fleqn,aspectratio=169,12pt]{beamer}

\usepackage{gaal}

\title{Circunferência}

%------------------------------------------------------------------------------%
\begin{document}

\addtitle

%------------------------------------------------------------------------------%
\section{Equação Canônica da Circunferência}

%------------------------------------------------------------------------------%
\begin{frame}{Equação Canônica da Circunferência}
  Equação \emph{canônica da circunferência}
  \pause
  \[
    x^2 + y^2 = r^2
  \]
  \pause
  com \(r > 0\)

  \pause
  \bigskip
  Circunferência centrada na origem de raio $r$
\end{frame}

%------------------------------------------------------------------------------%
\begin{frame}{Equação Matricial da Circunferência}
  Equação \emph{matricial da circunferência}
  \pause
  \[
    v^t I v = r^2
  \]
  com
  \[
    \pause
    v = \begin{bmatrix}
      x \\
      y
    \end{bmatrix}
    \qquad\qquad
    \pause
    I = \begin{bmatrix}
      1 & 0 \\
      0 & 1
    \end{bmatrix}
    \qquad\qquad
    \pause
    r > 0
  \]
\end{frame}

%------------------------------------------------------------------------------%
\begin{proofframe}{Pontos Relevantes da Circunferência}
\begin{columns}[t]
\column{0.4\textwidth}
  \onslide<+->Fazendo \(y=0\){}
  \begin{align*}
    \onslide<+->{x^2 + y^2 & = r^2 \\}
    \onslide<+->{x^2       & = r^2 \\}
    \onslide<+->{\abs{x}   & = r   \\}
    \onslide<+->{x         & = \pm r}
  \end{align*}

\column{0.4\textwidth}
  \onslide<+->{Fazendo \(x=0\)}
  \begin{align*}
    \onslide<+->{x^2 + y^2 & = r^2 \\}
    \onslide<+->{y^2       & = r^2 \\}
    \onslide<+->{\abs{y}   & = r   \\}
    \onslide<+->{y         & = \pm r}
  \end{align*}
\end{columns}
\end{proofframe}

%------------------------------------------------------------------------------%
\begin{frame}{Elementos Geométricos da Circunferência}
  \emph{Centro} da circunferência
  \[
    C = (0, 0)
  \]

  \pause
  \emph{Raio} da circunferência \(r\)

  \pause
  A circunferência cruza os eixos nos pontos
  \[
    (-r, 0) \qquad
    ( r, 0) \qquad
    (0, -r) \qquad
    (0,  r)
  \]
\end{frame}

%------------------------------------------------------------------------------%
\framedgraphic*{Gráfico da Circunferência}{definicao_circunferencia_1}{}
\framedgraphic*{Gráfico da Circunferência}{definicao_circunferencia_2}{}
\framedgraphic*{Gráfico da Circunferência}{definicao_circunferencia_3}{}
\framedgraphic*{Gráfico da Circunferência}{definicao_circunferencia_4}{}
\framedgraphic*{Gráfico da Circunferência}{definicao_circunferencia_5}{}

%------------------------------------------------------------------------------%
%------------------------------------------------------------------------------%
\begin{question}
  Dada a equação da circunferência
  \[
    x^2 + y^2 = 16
  \]
  \vspace{-1.5\baselineskip}
  \begin{enumerate}
    \item encontre seu raio e o centro
    \item escreva a equação matricial
    \item esboce a curva
  \end{enumerate}
\end{question}

%------------------------------------------------------------------------------%
\begin{resolution}{Raio de Centro}
  Comparando
  \[
    x^2 + y^2 = 16
  \]
  com a equação canônica
  \[
    x^2 + y^2 = r^2
  \]

  Identificamos

  \quad o centro:
  \(
    C = (0, 0)
  \)

  \quad o raio:
  \(
    r = \sqrt{16} = 4
  \)
\end{resolution}

%------------------------------------------------------------------------------%
\begin{resolution}{Equação Matricial}
  Equação Matricial
  \[
    v^t I v = r^2
  \]
  com
  \[
    I = \begin{bmatrix}
      1 & 0 \\
      0 & 1
    \end{bmatrix}
    \qquad\qquad
    r = 4
  \]
\end{resolution}

%------------------------------------------------------------------------------%
\begin{resolution}{Esboço}
\end{resolution}

%------------------------------------------------------------------------------%
\endinput


%------------------------------------------------------------------------------%
\section{Translação da Circunferência}

%------------------------------------------------------------------------------%
\begin{frame}{Translação da Circunferência}
  A \emph{translação} muda o centro, mas preserva a forma

  \pause
  \bigskip
  Utilizamos uma mudança de coordenadas para escrever a equação
\end{frame}

%------------------------------------------------------------------------------%
\framedgraphic*{Mudança de Coordenadas}{transformacao_translacao_1}{}
\framedgraphic*{Mudança de Coordenadas}{transformacao_translacao_2}{}
\framedgraphic*{Mudança de Coordenadas}{transformacao_translacao_3}{}
\framedgraphic*{Mudança de Coordenadas}{transformacao_translacao_4}{}
\framedgraphic*{Mudança de Coordenadas}{transformacao_translacao_5}{}

%------------------------------------------------------------------------------%
\begin{frame}{Equação da Circunferência Transladada}
  Novo centro \(C = (h, k)\)

  \pause
  Novas coordenadas
  \[
    \pause
    X = x - h
    \qquad
    \pause
    Y = y - k
  \]

  \pause
  Equação canônica nas novas coordenadas
  \[
    \pause
    X^2 + Y^2 = r^2
  \]
  \[
    \pause
    (x-h)^2 + (y-k)^2 = r^2
  \]
\end{frame}

%------------------------------------------------------------------------------%
\begin{proofframe}{Equação Matricial da Circunferência Transladada}
  Novo de centro
  \[
    \pause
    C
    \pause
    = \begin{bmatrix}
      h \\
      k
    \end{bmatrix}
  \]

  \pause
  \bigskip
  Novas coordenadas
  \[
    \pause
    V
    \pause
    = \begin{bmatrix}
      X \\
      Y
    \end{bmatrix}
    \pause
    = v - C
    \pause
    = \begin{bmatrix}
      x \\
      y
    \end{bmatrix}
    -
    \begin{bmatrix}
      h \\
      k
    \end{bmatrix}
    \pause
    = \begin{bmatrix}
      x - h \\
      y - k
    \end{bmatrix}
  \]
\end{proofframe}

%------------------------------------------------------------------------------%
\begin{proofframe}{Equação Matricial da Circunferência Transladada}
  \begin{align*}
    \onslide<+->{\emph{V^t I V}                    & = \emph{r^2}}   \\
    \onslide<+->{(v-C)^t I (v-C)                   & = r^2}          \\
    \onslide<+->{\emph{(x-h)^2 + (y-k)^2}          & = \emph{r^2}}   \\
    \onslide<+->{x^2 - 2hx + h^2 + y^2 - 2ky + k^2 & = r^2}          \\
    \onslide<+->{x^2 + y^2 - 2hx - 2ky + h^2 + k^2 & = r^2}
    \\
    \onslide<+->{\begin{bmatrix}
      x \\
      y
    \end{bmatrix}^t
    {\color<8->{magenta}\begin{bmatrix}
      1 & 0 \\
      0 & 1
    \end{bmatrix}}
    \begin{bmatrix}
      x \\
      y
    \end{bmatrix}
    +
    {\color<9->{purple}\begin{bmatrix}
      -2h \\
      -2k
    \end{bmatrix}^t}
    \begin{bmatrix}
      x \\
      y
    \end{bmatrix}
    + {\color<10->{violet}(h^2 + k^2 - r^2)} & = 0}
    \\
    \onslide<+->{v^t {\color<8->{magenta} I} v
    + {\color<9->{purple} B^t}v
    + {\color<10->{violet} f} & = 0}
  \end{align*}
\end{proofframe}

%------------------------------------------------------------------------------%
%------------------------------------------------------------------------------%
\begin{question}
  Dada a equação da circunferência
  \[
    (x-2)^2 + (y+3)^2 = 25
  \]
  \vspace{-1.5\baselineskip}
  \begin{enumerate}
    \item encontre seu raio e o centro
    \item escreva a equação matricial transladada
    \item esboce a curva
  \end{enumerate}
\end{question}

%------------------------------------------------------------------------------%
\begin{resolution}{Raio de Centro}
  Comparando
  \[
    (x-2)^2 + (y+3)^2 = 25
  \]
  com a equação canônica transladada
  \[
    (x-h)^2 + (y-k)^2 = 25
  \]

  Identificamos

  \quad o centro:
  \(
    C
    = (h, k)
    = (2, -3)
  \)

  \quad o raio:
  \(
    r
    = \sqrt{25}
    = 5
  \)
\end{resolution}

%------------------------------------------------------------------------------%
\begin{resolution}{Equação Matricial}
  Equação Matricial
  \[
    V^t I V = r^2
  \]
  com
  \[
    I = \begin{bmatrix}
      1 & 0 \\
      0 & 1
    \end{bmatrix}
  \]
  \[
    V
    = \begin{bmatrix}
      x - h  \\
      y - k
    \end{bmatrix}
    = \begin{bmatrix}[c]
      x - 2  \\
      y - (-3)
    \end{bmatrix}
    = \begin{bmatrix}
      x - 2  \\
      y + 3
    \end{bmatrix}
  \]

  \[
    r = 5
  \]
\end{resolution}

%------------------------------------------------------------------------------%
\begin{resolution}{Esboço}
\end{resolution}

%------------------------------------------------------------------------------%
\endinput

%------------------------------------------------------------------------------%
\begin{question}
  Determine o conjunto das soluções reais de
  \[
    (x-1)^2 + (y+2)^2 = 0
  \]
\end{question}

%------------------------------------------------------------------------------%
\begin{resolution}{Solução}
  Como os quadrados são não negativos
  \[
    (x-1)^2 \geq 0
    \qquad
    (y+2)^2 \geq 0
  \]
  A soma só é nula se
  \[
    x - 1 = 0
    \qquad
    y + 2 = 0
  \]
  Logo o conjunto de soluções é
  \[
    S = \{(1, -2)\}
  \]
  A equação representa uma circunferência degenerada em um ponto
\end{resolution}

%------------------------------------------------------------------------------%
\endinput


%%------------------------------------------------------------------------------%
%\section{Lista Mínima}
%
%%------------------------------------------------------------------------------%
%\begin{frame}{Lista Mínima}
%  \begin{enumerate}
%    \item
%  \end{enumerate}
%
%  \vfill
%  Atenção: A prova é baseada no livro, não nas apresentações
%\end{frame}

%------------------------------------------------------------------------------%
\end{document}
%------------------------------------------------------------------------------%
